\section{Purpose, scope and notation}\label{sec:intro}

\subsection{What the package does}

A component mounted on a vehicle, a machine or a package is shaken during its life. The acceleration
at its base is measured (a \emph{record}), or prescribed by a \emph{specification}, or described by a
\emph{mission}: a list of situations (road types, operating conditions) with the time spent in each.
An accelerated vibration test reproduces, on a shaker and in a shorter time $T$, the fatigue damage
that the life of duration $\Tl$ would cause. \code{vibsynth} computes the acceleration time history
$x(t)$ of such a test (the \emph{test profile}).

The fatigue damage is described by the Fatigue Damage Spectrum (FDS, Section~\ref{sec:fds}): the
damage that the vibration causes in a family of single-degree-of-freedom oscillators with different
natural frequencies. Two signals with the same FDS damage every lightly damped component in the same
way, whatever its natural frequency. Every synthesis method of the package (Parts~II and~III)
therefore starts from the target FDS and builds a signal with that FDS; the methods differ in what
they prescribe besides the FDS: the kurtosis of the signal, its RMS value, its power spectral density,
tonal components, or, in MMVSD, the kurtosis of the response of every oscillator. The \code{edit}
method of Part~IV takes the opposite route: it does not synthesise but keeps the damaging stretches of
the record itself, which makes it right for every S--N exponent, up to a limit that the record sets.

\subsection{How to read this manual}

The manual is written so that it can be followed without paper and pencil. Each symbol is defined at
the place where it first appears, and each derivation shows all intermediate steps. Results that are
used later are collected in boxes like this one:

\begin{result}[Example of a result box]
A boxed equation is a result that later sections refer to.
\end{result}

Three results are taken from the literature without derivation, because they are an empirical fit or
classical theorems: Dirlik's distribution of rainflow ranges, the central limit theorem and the
statistical properties of Welch's PSD estimate (Appendix~\ref{app:quoted}). They are stated precisely
and their role is explained. Everything else is derived here.

\subsection{Notation}

\begin{center}
\begin{tabular}{>{$}l<{$} p{0.72\textwidth}}
\toprule
t & time [s] \\
f,\ \omega = 2\pi f & frequency [Hz] and angular frequency [rad/s] \\
x(t) & base acceleration: the record or the test profile [m/s$^2$] \\
f_s,\ \Delta t = 1/f_s & sampling frequency and sampling interval \\
N,\ T = N\Delta t & number of samples and duration of a signal \\
\Tl & life represented by a target (the duration of use that the test must reproduce) \\
\E[\cdot] & expectation (average over all realisations of a random process) \\
\tavg{\cdot} & time average over the duration of a signal, $\tavg{g} = \frac1T\int_0^T g(t)\,\dd t$ \\
\sigma^2,\ \sigma & variance and standard deviation (RMS value of a zero-mean signal) \\
\kappa & kurtosis (Pearson definition, 3 for a Gaussian variable); $\kappa_z(\fn)$ is called
  \code{k\_z} in the program \\
G(f) & one-sided power spectral density (PSD) [(m/s$^2$)$^2$/Hz] \\
\fn,\ \wn = 2\pi\fn & natural frequency of an oscillator \\
Q,\ \zeta = 1/(2Q) & quality factor and damping ratio \\
z(t) & relative displacement of an oscillator (base minus mass) [m] \\
H(f) & frequency response from base acceleration to relative displacement \\
b,\ C,\ K & Basquin exponent, Basquin constant, stress per unit relative displacement \\
D(\fn) & Fatigue Damage Spectrum (FDS) \\
\mathrm{MRS}(\fn) & Maximum Response Spectrum \\
\Gamma(\cdot) & Euler's Gamma function, $\Gamma(s) = \int_0^\infty u^{s-1}e^{-u}\,\dd u$ \\
\bottomrule
\end{tabular}
\end{center}

The natural frequencies at which the FDS is evaluated form a grid $f_{\mathrm{n},1} <
f_{\mathrm{n},2} < \dots < f_{\mathrm{n},n}$ (the \emph{nodes}). By default the grid has 12 points per
octave: $f_{\mathrm{n},i} = f_{\min}\,2^{(i-1)/12}$, the last node being moved exactly to $f_{\max}$.
A linear grid with a constant step is also available. An index $i$ (or $j$) always runs over the
nodes, so that a quantity defined at every node, such as $D_i = D(f_{\mathrm{n},i})$, is a vector of
length $n$.

\subsection{Organisation}

Part~I (Sections~\ref{sec:signals}--\ref{sec:damage}) introduces random signals, the oscillator and
fatigue damage. Part~II (Sections~\ref{sec:synthesis}--\ref{sec:methods}) explains how a profile with
a prescribed FDS is designed, synthesised and corrected, and the methods of version~2.0. Part~III
(Sections~\ref{sec:mmvsdmodel}--\ref{sec:lines}) derives MMVSD and its treatment of
sine-on-random blocks. Section~\ref{sec:practice}
relates the theory to the settings of the program, and Section~\ref{sec:diagnostics} defines the
diagnostics and checks of the interface. Part~IV (Sections~\ref{sec:limit}--\ref{sec:nsm}, added in
version~2.2 after the practical sections so that the numbering of versions 2.0 and 2.1 is unchanged)
derives the compression limit of a record and the \code{edit} method, the kurtosis a compressed test
should carry, and the non-stationarity matrix behind the response kurtosis spectrum.
